\section{Problem Awareness} \label{sec:problem-awareness}

This chapter establishes the meta-requirements that motivate the design principle space developed in Chapter~\ref{sec:design-principles}, through a hermeneutic literature review. Section~\ref{ssec:review-methodology} presents the review methodology, the resulting literature corpus, and the three thematic directions into which the corpus is organised. The subsequent sections discuss the corpus along these three directions: human-\ac{AI} interaction with \ac{GenAI} systems (Section~\ref{ssec:hai-genai}), the attribution of \ac{ToM} to \ac{AI} systems (Section~\ref{ssec:tom-attributions}), and explanations and mental model formation (Section~\ref{ssec:explanations-mental-model-formation}). The meta-requirements derived from these sections are synthesised in Section~\ref{ssec:meta-requirements}.

\subsection{Review Methodology and Corpus} \label{ssec:review-methodology}

The hermeneutic literature review is an interpretive, iterative approach to reviewing literature \parencite{Boell2014}. Rather than exhaustively searching a predefined corpus, the process begins with an initial set of works which is iteratively expanded through search and acquisition cycles \parencite{Boell2014}. During the search new works are discovered (e.g. by tracking citations) and then acquired if they contribute relevant additions to the corpus. The iterative movement between individual works and the evolving whole is the defining characteristic of the hermeneutic circle \parencite{Boell2014}. The process terminates when saturation is reached, i.e., when new sources cease to produce interpretively relevant additions to the corpus.

Due to the nascent state of the field the hermeneutic approach was uniquely suited for the review. This was reinforced by the interdisciplinary nature of the research topic, as well as practical constraints of a master's thesis scope. The starting corpus consisted of the four works introduced and discussed in Section~\ref{ssec:related-works}: \textcite{Weisz2025}, \textcite{Jussupow2021}, \textcite{Colombatto2025}, and \textcite{Ying2024}. Their relevance to this thesis had already been established, making them a natural starting point for the iterative expansion of the corpus.

The corpus was expanded primarily by tracing the citations of the starting corpus forward and backward \parencite[cf.][]{Boell2014}. Forward citation search was conducted using \href{https://scholar.google.com/}{Google Scholar}\footnote{https://scholar.google.com}, but yielded limited results due to the recency of the starting corpus, which is itself a reflection of the nascent state of the field. Backward citation search proved more productive: tracing the references of \textcite{Ying2024} led to \textcite{Shin2021}, whose references in turn led to \textcite{Arrieta2020} and \textcite{Miller2019}. \textcite{Amershi2019} was identified through \textcite{Weisz2025}'s related conference publication on the same design principle framework \parencite{Weisz2024}, which references it directly. \textcite{Hoffman2023} was identified through topical search rather than citation tracking.

Search was bounded thematically to topics relevant to the meta-requirements of this thesis, as established in the previous chapters. Additionally, a temporal boundary was set to 2017 onward, marking the introduction of the transformer architecture \parencite{Vaswani2017} as the technical inflection point that enabled modern \ac{GenAI} systems \parencite{Schneider2024}. Seminal works predating this boundary were included where they provide theoretical foundations that have not been superseded. Saturation was considered reached when new sources were found to either duplicate positions already represented in the corpus or focus on implementation details rather than conceptual or empirical contributions relevant to the meta-requirements.

The resulting corpus of nine works is organised into three thematic directions, which structure the remainder of this chapter:

\begin{itemize}
    \item \enquote{Human-\ac{AI} interaction with \ac{GenAI}}, covered in Section~\ref{ssec:hai-genai} and centred on \textcite{Weisz2025} and \textcite{Amershi2019}.
    \item \enquote{Theory of Mind attributions to \ac{AI} systems}, covered in Section~\ref{ssec:tom-attributions} and centred on \textcite{Jussupow2021}, \textcite{Colombatto2025}, and \textcite{Ying2024}.
    \item \enquote{Explanations and mental model formation}, covered in Section~\ref{ssec:explanations-mental-model-formation} and centred on \textcite{Miller2019}, \textcite{Shin2021}, \textcite{Hoffman2023}, and \textcite{Arrieta2020}.
\end{itemize}

% TODO: If there is time, consider finding a second source that was excluded. The goal would be a "standard work" that is widely cited in the field but was excluded
Saturation was reached with the inclusion of \textcite{Hoffman2023}, which provided new insights into the measurement of mental models and their relationship to explanations. Additional works \parencite[e.g.,][]{Baker2017} concerned themselves with the design of \ac{AI} systems with \ac{ToM} capabilities which is a related but distinct topic that was not included in the corpus as it did not contribute to the meta-requirements of this thesis.

\subsection{Human-AI Interaction with GenAI} \label{ssec:hai-genai}



\subsection{Theory of Mind Attributions to AI Systems} \label{ssec:tom-attributions}



\subsection{Explanations and Mental Model Formation} \label{ssec:explanations-mental-model-formation}

% consider \parencite{Halpern2005,Kashima1998} for the role of user beliefs in shaping perception and interaction with AI systems, and how this relates to ToM attribution

\subsection{Meta-Requirements} \label{ssec:meta-requirements}
